% preliminares/prologo.tex
% Fuente: Prólogo.docx
\addchap{Prólogo}
\addtocontents{toc}{\protect\autorentoc{Mabel Giammatteo y Laura Ferrari}}

La gramática, surgida en la Antigüedad clásica, venerada y vilipendiada a lo
largo de su desarrollo en Occidente, es actualmente considerada «la estructura
del lenguaje». Según \textcite{saussure1945}, este sistema es el encargado de
aparear determinadas porciones de la masa amorfa de nuestros pensamientos con
determinadas porciones del sonido para conformar signos lingüísticos que hagan
comunicables nuestros pensamientos, deseos y sentimientos. Por su complejidad,
esta estructura ha sido comparada con el código genético, ya que, así como este,
mediante la combinación de una mínima cantidad de elementos, puede generar la
inmensa variedad de lo viviente, la gramática nos permite «combinar un número
relativamente pequeño de unidades y obtener un número potencialmente infinito de
resultados» \parencite[18]{bosque2009}. Dado, entonces, que tanto para precisar
nuestras ideas como para interactuar con otras personas nos valemos del
lenguaje, podemos afirmar que nuestro desarrollo cognitivo y nuestra vida social
están bajo el dominio de la gramática.

Por un lado, la gramática se nos revela como un cuerpo de doctrina atesorado por
siglos de estudios sobre el lenguaje y las distintas lenguas y, por el otro,
constituye un conocimiento vivo, según iremos viendo, internalizado en la mente
de cada hablante \parencite{chomsky1957} e indispensable para la interacción
social. Este manual surge en la intersección entre el saber teórico y la
necesidad de hacer consciente el conocimiento sobre cómo funcionan las
estructuras que hacen posible el lenguaje. Es fruto de un proyecto compartido
nacido en el seno de la cátedra B de Gramática de la Facultad de Filosofía y
Letras de la Universidad de Buenos Aires. Su objetivo ha sido ofrecer las
nociones y herramientas básicas para analizar la lengua, pero, sobre todo,
propiciar un acercamiento reflexivo, que permita comprender cuáles son sus
unidades y cómo se combinan para construir significados relevantes.

La obra, que recoge la experiencia docente y el trabajo de investigación del
equipo, surgió inicialmente como una simple guía de ejercicios para los y las
estudiantes de nuestros cursos de primer año de la Facultad y fue poco a poco
creciendo con la incorporación de las nociones teóricas que fundamentaban las
distintas prácticas solicitadas. Asimismo, el núcleo inicial de temas abordados
también se fue ampliando hasta convertirse en lo que hoy presentamos: un manual
completo teórico-práctico, que refleja las teorías actuales sobre el sistema y
las aplica al estudio de la lengua de todos los días.

Se trata de un texto colectivo, con un espíritu común, pero cuyos capítulos, que
pueden ser aprovechados en forma independiente, conservan la impronta que les
han asignado quienes los han producido. La idea central ha sido proveer un
material actualizado, a la vez que accesible, que pueda servir de guía o hilo
conductor para cualquier curso de gramática del español en las etapas iniciales
del nivel superior, ya sea en la universidad o en profesorados de Lengua y
Literatura, pero que también puede ser de utilidad en cursos superiores de grado
y posgrado y en capacitaciones docentes. Asimismo, también puede servir de base
de comparación para quienes estudian alguna lengua extranjera, y como texto de
referencia para profesionales que trabajan con el lenguaje, como periodistas,
especialistas en traducción, corrección de textos o edición, y para quienes estén
interesados en profundizar sus conocimientos sobre los distintos aspectos de la
gramática del español.

El libro se inicia con un capítulo introductorio que presenta diferentes teorías
lingüísticas y su incidencia en la enseñanza de la gramática, seguido de otro
dedicado a distinguir entre nociones básicas, como las de gramaticalidad,
corrección y aceptabilidad. Los capítulos siguientes conforman un trayecto que
aborda las principales cuestiones de la gramática del español, desde los sonidos
y su caracterización, pasando por la estructura de la palabra y sus procesos de
formación, tanto los tradicionales como otros más informales y novedosos, hasta
la actitud comunicativa del hablante y su reflejo en la modalidad de las
oraciones. La parte central de la obra considera la estructura oracional desde
una perspectiva que integra tres niveles: el nivel interior o del sintagma
verbal, relacionado con la predicación; el nivel central, vinculado con la
flexión, donde se establece la concordancia entre el sujeto y el predicado; y el
nivel externo o periferia oracional, en el que se manifiesta la fuerza ilocutiva
del hablante. Se tratan también cuestiones relacionadas con aspectos
problemáticos de la gramática del español, como las construcciones pronominales,
las perífrasis verbales o los predicativos, y los nuevos planteos para el
tratamiento de las estructuras interoracionales, con un capítulo dedicado a las
oraciones coordinadas y otro a las subordinadas. La obra incluye un capítulo
final dedicado a la variación, en el que se discuten las nociones de lengua y
dialecto y se presentan las principales características del español del Río de
la Plata.

Las nociones teóricas se abordan desde un enfoque léxico-sintáctico que no solo
tiene en cuenta las estructuras de la lengua, sino también los significados que
resultan de sus combinaciones y las interpretaciones que producen en los textos
en los que se insertan. Así, lo más destacable de la obra es que no se limita a
la descripción de las formas lingüísticas ni a su explicación en relación con
principios gramaticales generales, sino que las estructuras se conectan con el
uso en diferentes tipos de textos, tanto literarios como cotidianos (noticias
periodísticas, afiches callejeros, memes de internet, canciones y proverbios,
entre otros géneros que la obra aborda). Cada capítulo contiene numerosos
ejemplos que dan cuenta de las estructuras estudiadas y proporciona abundantes y
variadas actividades que permiten poner en práctica los contenidos. La obra
también incorpora un solucionario con las claves de resolución de los
ejercicios, y se completa con un siglario y un glosario que hacen más accesibles
su empleo y su consulta.

Hasta aquí, nuestro trabajo se ha concentrado en el diseño de la estructura del
libro, en la supervisión de la escritura de los capítulos y en la revisión
final. A partir de este momento, con el texto ya concluido, solo nos queda
confiar en que sea de utilidad para quienes deciden aventurarse por los caminos
del estudio de la gramática y también en que despierte el interés para seguir
leyendo. Para ello, los distintos capítulos proporcionan una lista de textos de
referencia que permitirán ampliar y profundizar los conocimientos.

\bigskip
\begin{flushright}
  \itshape Las coordinadoras\\
  \upshape Buenos Aires, 9 de octubre de 2026
\end{flushright}
